\section{Inventario de sistemas}

El inventario de sistemas constituye el punto de partida para comprender el ecosistema de información de Franco's S.A.C. en su estado actual (AS-IS). Su propósito es identificar con qué herramientas registra, procesa y comparte información la organización, quiénes las utilizan, sobre qué infraestructura operan y en qué medida se relacionan entre sí. Para ello, el levantamiento se organizó en dos niveles: los sistemas y registros que soportan la operación, y la infraestructura tecnológica sobre la cual estos se ejecutan.

\subsection{Sistemas y registros}

Los elementos identificados se agruparon en cuatro tipos según su naturaleza: sistema empresarial, registros manuales, herramientas ofimáticas y medios de comunicación. Esta clasificación permite distinguir entre las soluciones formalmente adquiridas por la empresa y aquellas que el personal emplea de manera informal para suplir la ausencia de un sistema integrado.

\begin{table}[H]
\centering
\caption{Inventario de sistemas y registros de Franco's S.A.C. (AS-IS)}
\label{tab:inventario-sistemas}
\small
\renewcommand{\arraystretch}{1.3}
\begin{tabularx}{\textwidth}{|>{\raggedright\arraybackslash}p{2.0cm}|>{\raggedright\arraybackslash}p{2.0cm}|X|>{\raggedright\arraybackslash}p{2.3cm}|>{\raggedright\arraybackslash}p{2.6cm}|}
\hline
\textbf{Tipo} & \textbf{Sistema} & \textbf{Función} & \textbf{Usuarios} & \textbf{Plataforma} \\
\hline
Sistema empresarial & SICO ERP & Emisión de comprobantes de pago (boletas y facturas); sus demás módulos no se utilizan. & Personal de ventas y facturación & Instalación local en una única PC, sin servidor ni nube \\
\hline
Sistema empresarial & SICO Móvil & Aplicación móvil complementaria a SICO ERP; fue adquirida, pero no se encuentra en uso. & Sin usuarios & No implementada \\
\hline
Registro manual & Formatos físicos & Registro de ventas, distribución e inventario, sin un formato estandarizado. & Personal operativo de ventas, inventario y distribución & Papel, en planta y tienda \\
\hline
Ofimática & Excel & Registro y cálculo de información de forma informal. & Personal de la empresa & PC local con Windows \\
\hline
Ofimática & Word & Elaboración de documentos de forma informal. & Personal de la empresa & PC local con Windows \\
\hline
Comunicación & Correo electrónico & Comunicación formal interna y externa. & Personal de la empresa & Servicio web, vía ISP local \\
\hline
Comunicación & WhatsApp & Coordinación informal entre el personal. & Personal de la empresa & Servicio web \\
\hline
\end{tabularx}
\end{table}

Como se observa en la Tabla~\ref{tab:inventario-sistemas}, SICO ERP es el único sistema empresarial en operación; sin embargo, su aprovechamiento es mínimo, pues se limita a la facturación y deja sin uso el resto de sus funcionalidades. Esta subutilización se extiende a SICO Móvil, que pese a formar parte de la adquisición no ha sido implementado. En consecuencia, la información operativa de ventas, inventario y distribución se registra en papel o en archivos de Excel y Word elaborados sin un criterio común, mientras que la coordinación diaria descansa en el correo electrónico y en WhatsApp. Ninguno de estos elementos intercambia información con otro: SICO ERP solo se vincula con SICO Móvil por pertenecer a la misma adquisición, y los registros manuales y ofimáticos permanecen aislados.

\subsection{Infraestructura tecnológica}

La infraestructura que soporta estos sistemas es reducida y de tipo doméstico. La empresa opera sobre una única red local cableada, sin segmentación, y no cuenta con servidor propio, servicios en la nube ni herramientas de monitoreo. La Tabla~\ref{tab:inventario-infraestructura} describe los componentes identificados.

\begin{table}[H]
\centering
\caption{Inventario de infraestructura tecnológica de Franco's S.A.C. (AS-IS)}
\label{tab:inventario-infraestructura}
\small
\renewcommand{\arraystretch}{1.3}
\begin{tabularx}{\textwidth}{|>{\raggedright\arraybackslash}p{2.6cm}|X|>{\centering\arraybackslash}p{1.4cm}|>{\raggedright\arraybackslash}p{2.3cm}|}
\hline
\textbf{Componente} & \textbf{Función} & \textbf{Cantidad} & \textbf{Ubicación} \\
\hline
PC con Windows & Estaciones de trabajo del personal; solo una de ellas ejecuta SICO ERP. & 4 & Oficina / planta \\
\hline
Impresora & Impresión de documentos. & 1 & Oficina \\
\hline
Switch de red & Interconexión de los equipos en la red local cableada. & 1 & Oficina / planta \\
\hline
Router (módem del ISP) & Salida a internet de la red local. & 1 & Oficina / planta \\
\hline
Servicio de internet & Conexión externa provista por un ISP local. & 1 & Proveedor externo \\
\hline
\end{tabularx}
\end{table}

La configuración descrita implica que toda la operación de facturación depende de un solo equipo. Al no existir un servidor ni un mecanismo de respaldo, cualquier falla en esa PC compromete la continuidad de la emisión de comprobantes y pone en riesgo la información almacenada en la base de datos de SICO ERP.

En síntesis, el inventario evidencia un ecosistema de información mínimo y desarticulado, caracterizado por tres condiciones: un sistema empresarial subutilizado y dependiente de un único equipo sin respaldo; datos operativos dispersos en papel y archivos sueltos, sin una base de datos central; y la ausencia total de integración entre las herramientas existentes. Estas condiciones se analizan con mayor detalle en las secciones siguientes.
